% Closed foundations for Kamp III.2, Lemma 8; full elementary induction remains separate.

\section{Finite-family induction and the complete boundary decomposition}
\label{sec:annex-p09-kamp-iii2-lemma8-foundations}

\begin{annexitem}{annex:p9.kamp.family.rank}{thm:p09-kamp-family-rank}{The exact source comparison and the residual-family bound}
\begin{proof}
All families below are finite sets of formulas with the positive integer
lengths defined in the body. Write \(m_F=M(F)\), \(m_G=M(G)\).
Suppose first that \(G\prec F\), witnessed by \(m\).
The strict inequality of counts implies that \(F\) has at least
one member of length \(m\). Hence \(F\ne\varnothing\),
\(m\geq1\), and \(m\leq m_F\). Every member of \(G\)
has length at most \(m\), so \(m_G\leq m\), including
the empty-family convention. It follows that \(m_G\leq m_F\).
If these maxima differ, then \(m_G<m_F\), as required.
If they are equal, the two inequalities force \(m=m_G=m_F\),
and the source's strict count inequality is exactly \(c(G)<c(F)\).

Conversely, if \(m_G<m_F\), choose \(m=m_F\).
Every member of \(G\) has length below \(m\), whereas
\(F\) has a member of length \(m\). Thus the relevant
counts are zero and positive, which witnesses \(G\prec F\).
If the maxima agree and \(c(G)<c(F)\), choose their common
maximum for \(m\). Every member of \(G\) meets the length
bound, and the required count inequality holds. The empty family
has no predecessor: its counts are all zero. This proves the
equivalence in every case.

For well-foundedness, consider any nonempty collection of finite families.
Among their natural-number maxima choose the least value \(m\).
Among the families with that maximum choose the least value of
\(c\), and choose any family attaining it. The equivalence just
proved shows that no member of the collection precedes this family.
Thus every nonempty collection has a \(\prec\)-minimal member,
which proves well-foundedness directly. One can equivalently observe
that in a descending chain the maxima cannot decrease infinitely,
and after their final decrease the nonnegative counts must decrease
at every step, which is impossible.

If a maximal-length member \(\phi\) is deleted, either another
member of that same length remains and its count falls by exactly one,
or the maximum itself decreases. This includes deleting the only member,
which gives the empty family. Hence deletion strictly decreases \(\prec\).

Now let \(m=M(F)>1\) and delete one maximal member. No permitted
replacement has length above its original member's length, so none
has length above \(m\). A remaining original member shorter than
\(m\) supplies no length-\(m\) output. Each remaining original
member of length \(m\) supplies at most one such output. There
are only \(c(F)-1\) of those members after deletion. Consequently
the union of all replacement sets has at most \(c(F)-1\) members
of length \(m\). Any coincidences between output formulas reduce
the cardinality of this union and cannot increase it. Retained or
new length-one formulas also contribute no member of length \(m\),
since \(m>1\). Thus the final family has maximum at most \(m\)
and fewer length-\(m\) members than \(F\), witnessing \(G\prec F\).
No bound on the number of shorter outputs was needed.

For a cut residual from Lemma~7, the original word has length
\(2n+1\). Its point cuts and gap cuts were computed explicitly
in \cref{thm:p09-kamp-cut-factorization}: all residual lengths
are at most the original and each orientation has just one possible
equal-length cut. Every subset alternative includes that index at
most once. Therefore each such residual satisfies the replacement
hypothesis exactly, including an empty residual. The source's argument
on p.~91 uses this bound, not an enlarged word length
\cite[ch.~III, pp.~83, 91]{Kamp1968}.
\end{proof}
\end{annexitem}

\begin{annexitem}{annex:p9.kamp.family.cut}{thm:p09-kamp-family-cut}{Every tuple of cut choices, both orientations and the exact residual count}
\begin{proof}
For each original word the index set has \(2n_j-1\) members,
so its power set has \(2^{2n_j-1}\) members. Multiplying these
finite cardinalities gives the stated cardinality of the Cartesian
product for \(k>0\). For \(k=0\) there is exactly one
choice, the function on the empty index set, represented by the
empty tuple. Thus the complete enumeration is finite in every case.

Fix a linear structure, its predicate interpretation, and an assignment
to every free parameter, with \(a,b\) the values of \(x,y\).
Suppose first that the original left-interior formula is true.
Choose its outer witness \(d\), so \(a<d<b\) and
\(\psi_1\) holds. For each \(j=1,\ldots,k\) define
\[
             J_j=\{c\in\mathcal C_{n_j}:R_{j,c}(d,b)\text{ holds}\}.
\]
These are all the fixed-side truth subsets at the same witness.
Their tuple belongs to the enumerated product and its aggregate
\(S_{\mathbf J}^L(d,b)\) holds, since every selected conjunct
is true. An empty conjunction is true by convention.

For any \(t\) with \(a<t<d\), if \(\psi_0\) is true
with \(v=t\), the new matrix holds. Otherwise the original
matrix gives \(\rho_j(t,b)\) for at least one original
index \(j\). Apply the cut factorization to that word at
\(t<d<b\). It supplies some cut \(c\) with
\(L_{j,c}(t,d)\land R_{j,c}(d,b)\).
The second conjunct puts \(c\) in the already selected
\(J_j\), and the first supplies one residual disjunct.
Thus the new universal matrix holds for every tested \(t\),
and the same \(d\) witnesses the tuple's output alternative.
If \(k=0\), the original disjunction is false, so the case
where \(\psi_0\) is false never occurs; the same reasoning
gives the sole empty-tuple alternative.

Conversely, suppose an output alternative is true with tuple
\(\mathbf J\) and witness \(d\). Its fixed-side aggregate
gives \(R_{j,c}(d,b)\) for every selected pair \((j,c)\).
For any \(a<t<d\), either \(\psi_0\) holds or the
residual disjunction supplies some selected pair with
\(L_{j,c}(t,d)\). In the latter case combine that condition
with \(R_{j,c}(d,b)\) and the reverse cut factorization to
obtain \(\rho_j(t,b)\). Hence the original matrix holds
at every point. A false empty residual cannot supply a pair,
so the converse also covers an empty tuple or all-empty subsets.
The same witness supplies the unchanged guard and \(\psi_1\).

For the right-interior equivalence, suppose its original formula
holds with witness \(d\), where \(a<d<b\). For each word
select
\[
             J_j=\{c\in\mathcal C_{n_j}:L_{j,c}(a,d)\text{ holds}\}.
\]
Every selected prefix is true, giving \(S_{\mathbf J}^R(a,d)\).
Take any \(d<t<b\). If \(\psi_0\) is true there, the
new matrix holds. Otherwise choose an original \(\rho_j(a,t)\)
supplied by its original disjunction. Factor this word at
\(a<d<t\), obtaining
\(L_{j,c}(a,d)\land R_{j,c}(d,t)\).
The first puts the index in the fixed subset \(J_j\), and
the second gives the corresponding residual. Thus the new matrix
holds everywhere and its tuple alternative is true.

In the reverse right direction, a true alternative supplies
\(L_{j,c}(a,d)\) for each selected pair. At each tested
\(d<t<b\), a false \(\psi_0\) forces some residual
\(R_{j,c}(d,t)\). Combining the two cut factors yields
\(\rho_j(a,t)\), as required by the original matrix.
If the aggregate residual is empty this case is impossible, so
the same converse includes every empty configuration. This proves
both orientations without omitting any tuple or any direction.

No fixed subset depends on the universally varied point. A selected
word or cut may nevertheless differ from one point to the next;
the residual disjunction expresses precisely that possibility.
The original \(\psi_0,\psi_1\), their additional parameters,
and all their quantifier scopes remain unchanged. In particular
free occurrences of \(v\) in \(\psi_1\) remain outside the
inner universal quantifier. All newly constructed patterns have
fresh internal bound variables and their displayed free endpoints.

For each original \(j\), its residual cuts have length at
most \(2n_j+1\), with at most one candidate of that length
on the varying side. Its selected subset cannot increase this
number. Taking the union of residual formulas can identify equal
outputs and therefore only reduce the number of maximal-length
members. After deletion of a distinguished maximal member,
\cref{thm:p09-kamp-family-rank} applies to the resulting
residual family, provided the list enumerates each replaced family
member once and the other retained members obey that lemma's
hypotheses. An arbitrary auxiliary \(\psi_0\), although allowed
in the formula equivalences, is not automatically a finite family
of the required At/Bp form. Empty residuals and added length-one
formulas meet the explicit numerical boundary conditions.

The fixed-side conjunctions are separate components of every
alternative and have not been dropped in this count. Their
truth was indispensable in both converses. The construction
therefore supplies precisely the repeated finite choice step
underlying source pp.~90--91, while the later assembly of these
components into the required interval representations is a
separate obligation \cite[ch.~III, pp.~80--83, 90--91]{Kamp1968}.
\end{proof}
\end{annexitem}

\begin{annexitem}{annex:p9.kamp.family.normalization}{thm:p09-kamp-family-normalization}{The point-only base and complete restoration of fresh labels}
\begin{proof}
Suppose first that \(F\subseteq\operatorname{At}(z)\).
Every member belongs to \(\operatorname{E}(z)\), and that class
is closed under finite nonempty disjunctions. When the family is empty,
its selected falsehood \(F^K_z\) also belongs to that class.
Thus \(D_F\) has an explicit \(\operatorname{E}\) construction
in every case. Its renamed version \(D_F(w)\) and the truth
labels form the elementary word
\(W=(T^K_w,D_F(w),T^K_w)\).

For ordered endpoints \(a<b\), a realization of this one-interval
word is just the pair \((a,b)\). Its endpoint tests are true,
and its sole interval test is exactly
\(\forall t\,(a<t<b\to D_F(t))\). Hence
\(\mathcal B_p(W;x,y)\) has an \(\operatorname{Ep}(x,y)\)
presentation and agrees with \(\mathcal U_F\) under \(x<y\)
over all linear orders, by
\cref{thm:p09-kamp-elementary-intervals}.
Restricting that identity to \(\mathsf{Com}\) gives the required
property \(I(x,y)\). For the empty family the interval label
is false, so the formula holds precisely when the open interval
is empty. No additional point is required or supplied.

For the fresh-label reduction, enumerate the finitely many members of
\(F\) and every predicate occurrence in its basic members. Choose
finitely many distinct fresh predicate letters, all absent from \(F\)
and from the source truth/falsehood letter \(P_0\).
For each whole \(\operatorname{At}(z)\) member \(\theta\),
put one fresh positive atom \(Q_\theta(z)\) in its place.
For each basic member, keep its exact grammar, variable binding and
occurrence order, but replace each individual predicate occurrence by
its own fresh letter. The resulting basic formula has the original
number of point and interval positions and the same word length.
An \(\operatorname{At}\) member and its fresh atom both have
the declared source length one, regardless of the original number
of literals. This also treats the source's particular tautology:
one fresh atom represents that whole formula, rather than turning
its two repeated predicate occurrences into a formula outside the
declared \(\operatorname{At}\) grammar.

Every resulting formula contains a fresh letter used by no other
member, so distinct input members give distinct output members.
Thus the construction gives a bijection between the finite families
and preserves all their length counts. Each fresh letter occurs
exactly once in the new family, and no old letters remain there.
In particular \(M\), \(c\), and all predecessor comparisons
are preserved.

Restore the fresh atoms by a finite sequence of proper substitutions.
For \(Q_\theta\), use the original \(\theta(w)\), which
belongs to \(\operatorname{E}(w)\). For the fresh letter at
a basic predicate occurrence originally named \(P_j\), use
\(P_j(w)\), also in \(\operatorname{E}(w)\). Choose the
original enumeration to fix the order of these substitutions.
No replacement contains a fresh target letter, since all targets were
chosen outside the original vocabulary. Hence a later step cannot
alter a replacement inserted earlier. The operation restores exactly
the original atoms and complete \(\operatorname{At}\) members;
any bound-variable renaming is the allowed proper-substitution variant.
Order atoms, guards, endpoint variables and the universal quantifier
are retained. Therefore the resulting formula is
\(\mathcal U_F\), up to that bound-name variation.

By the proved substitution closure of property \(I\),
\cref{thm:p09-kamp-interval-substitution}, every step transports
an elementary interval representation. This proves the conditional
reduction from \(G\) back to \(F\). It is not a claim that
\(\mathcal U_G\) already has the required representation.

Finally suppose \(M(G)>1\), and add a fresh positive atom \(Q(z)\).
It has length one and cannot duplicate an existing member, because
its letter is fresh. It changes neither the maximum length nor
the number of members at that maximum. The comparison lemma therefore
shows that \(G\) and \(G^Q=G\cup\{Q(z)\}\) have exactly
the same predecessors. In the universal disjunction, substitute
the falsehood \(F^K_w\) properly for \(Q\). Nothing in
\(\mathcal U_G\) is affected, and the new matrix is
\(D_G(z)\lor F^K_z\), equivalent to \(D_G(z)\).
The same substitution closure transports any supplied
\(I(x,y)\) representation for \(\mathcal U_{G^Q}\) to one
for the original \(\mathcal U_G\).

Falsehood is essential at this step. For example, let
\(G=\{\phi(x,z)\}\) with a one-interval basic word,
take the order \(0<1<2\), and interpret every letter of
\(\phi\) by the empty set. At \(x=0,y=2\),
\(\mathcal U_G\) is false at its interior point 1.
Substituting a tautology for the added disjunct instead makes the
universal matrix true at every point and gives a true formula.
That substitution cannot recover the original condition. The active
falsehood formula above directly supplies the required correction
to the displayed fresh-disjunct step on source p.~84
\cite[ch.~III, pp.~83--84]{Kamp1968}.
\end{proof}
\end{annexitem}

\begin{annexitem}{annex:p9.kamp.fixed.endpoint.limits}{thm:p09-kamp-fixed-endpoint-limits}{The complete fixed-endpoint accumulation constructions}
\begin{proof}
Fix a linear structure and ordered endpoints \(a<b\).
The \(\operatorname{Lp},\operatorname{Lf}\) clauses in
\cref{def:p09-kamp-limiting-classes} and the basic realization
equivalence give the following exact meanings for the new labels.
For \(n=1\), \(\lambda_W^+(a)\) says that an
\(A_1\)-point occurs strictly between \(a\) and every point
greater than \(a\). For \(n>1\), it says that the entire
suffix word has a realization strictly inside each such interval.
The label \(\lambda_W^-(b)\) has the corresponding meaning
for the point \(A_0\), or the entire prefix word, strictly
inside every interval ending at \(b\). This follows directly
from the displayed limiting-class quantifiers; it makes no claim
that a one-letter word is a positive-stage basic pattern.

Suppose \(C_W^+(a)\) holds. Its instance at \(v=b\)
supplies a realization of the full word from \(a\) to some
\(q<b\). Write it as
\(a=a_0<a_1<\cdots<a_n=q\).
Then \(A_0(a)\) holds. Put \(u=a_1\); at \(n=1\)
this is the final point \(q\). In every case \(a<u<b\)
and \((a,u)\) satisfies \(B_1\) throughout.

For any \(v>a\), the same defining condition supplies a
full realization with final point below \(v\). Deleting its
initial point \(a\) leaves the suffix word strictly inside
\((a,v)\) when \(n>1\). At \(n=1\), it leaves a
single \(A_1\)-point there. Thus every required interval
satisfies the exact condition for \(\lambda_W^+(a)\).
Together these assertions give
\[
 A_0(a)\land\lambda_W^+(a)\land
       \exists u\,(a<u<b\land\forall t\,(a<t<u\to B_1(t))).
 \tag{1}
\]

Conversely, suppose (1) holds, and choose its witness \(u\).
For any \(v>a\), let \(d\) be the smaller of \(u,v\).
Then \(a<d\), and \(\lambda_W^+(a)\) supplies either
one \(A_1\)-point \(q\) in \((a,d)\), if \(n=1\),
or a suffix realization
\(a<a_1<\cdots<a_n<d\), if \(n>1\).
In the first case prepend \(a\) to \(q\); in the second
prepend \(a\) to the suffix sequence. The new first gap is
contained in \((a,u)\), so it satisfies \(B_1\).
The initial point satisfies \(A_0\), and every suffix point
and gap retains its already supplied label. The resulting sequence
realizes the whole original word with final point strictly below
\(d\leq v\). This is the required witness for
\(C_W^+(a)\). Since \(v>a\) was arbitrary, the reverse
implication is proved.

The two-interval word \(V_W^+\) has its first point condition
\(A_0\land\lambda_W^+\), its first gap condition \(B_1\),
and true labels at the intermediate point, second gap and final
point. Its realization from \(a\) to \(b\) therefore means
exactly (1). The labels have explicit \(\operatorname{E}\)
constructions, so \cref{thm:p09-kamp-elementary-intervals}
gives an \(\operatorname{Ep}(x,y)\) presentation. This
proves the first guarded \(\mathsf{Lin}\) identity.

For the other orientation, suppose \(C_W^-(b)\) holds.
Its instance at \(v=a\) supplies a full realization
\(a<q=a_0<\cdots<a_n=b\).
Its penultimate point \(u=a_{n-1}\) satisfies \(a<u<b\);
at \(n=1\) it is the initial point \(q\). Its final
gap \((u,b)\) satisfies \(B_n\), and \(A_n(b)\) holds.
For each \(v<b\), the defining condition supplies a full
realization whose initial point lies strictly above \(v\).
Delete its last point \(b\). For \(n>1\), this leaves
a realization of the entire prefix strictly inside \((v,b)\);
for \(n=1\), it leaves a single \(A_0\)-point there.
Thus \(\lambda_W^-(b)\) holds, and we have
\[
 A_n(b)\land\lambda_W^-(b)\land
       \exists u\,(a<u<b\land\forall t\,(u<t<b\to B_n(t))).
 \tag{2}
\]

Conversely, choose a witness \(u\) for (2). For any \(v<b\),
let \(d\) be the larger of \(u,v\), so \(d<b\).
The label \(\lambda_W^-(b)\) supplies a single
\(A_0\)-point in \((d,b)\) when \(n=1\), or an
entire prefix realization strictly inside \((d,b)\) when
\(n>1\). Append \(b\). The final gap is contained in
\((u,b)\) and satisfies \(B_n\); the last point satisfies
\(A_n\), and all prefix conditions are preserved. The full
word is realized from a point strictly above \(d\geq v\)
to \(b\), which proves the required instance of
\(C_W^-(b)\). This works for every \(v<b\).

The word \(V_W^-\) represents exactly (2): its first point,
first gap and intermediate point have true labels, while its
last gap and last point require \(B_n\) and
\(A_n\land\lambda_W^-\). It therefore has an elementary
past presentation and proves the second guarded identity.
Restricting either identity to complete orders gives property
\(I(x,y)\), without an appeal to the unfinished Lemma~8.

The strict guard \(a<b\) is essential to the displayed
representations: it supplied the instances used to extract actual
realizations. At adjacent endpoints each two-interval representative
is false, and the corresponding accumulation condition fails at
the other endpoint's quantified instance. Empty gaps elsewhere in
the witnesses retain their ordinary vacuous interval conditions.
The proof uses only minimum or maximum of two already given points,
not order completeness. It covers both \(n=1\) and every
positive-stage suffix or prefix without reducing their required
word length \cite[ch.~III, pp.~89--90]{Kamp1968}.
\end{proof}
\end{annexitem}

\begin{annexitem}{annex:p9.kamp.eleven.boundaries}{thm:p09-kamp-eleven-boundaries}{Existence, uniqueness and every boundary configuration}
\begin{proof}
Fix a complete linear order, ordered points \(a<b\), and the
truth set \(R=R_{a,b}\subseteq(a,b)\). The following order
argument is independent of the internal formula defining \(H\),
and therefore applies to both distinguished-pattern orientations.

If \(R\) is empty, case 1 holds and no nonempty case does.
Suppose \(R\ne\varnothing\). The point \(a\) is a lower
bound and \(b\) an upper bound, so completeness supplies
\(l=\inf R\) and \(h=\sup R\) in the domain.
For each \(t\in R\),
\[
                          a\leq l\leq t\leq h\leq b.
\]
In fact \(l<b\), since any \(t\in R\) has \(t<b\)
and \(l\leq t\). Similarly \(a<h\).

If \(l=h\), every member of \(R\) equals this common point,
by the displayed inequalities and antisymmetry of the associated
nonstrict order. Since the set is nonempty, it is exactly the
singleton of that point. Conversely, a singleton \(\{c\}\)
has infimum and supremum both equal to \(c\). Its membership
in \((a,b)\) gives \(a<c<b\), so this is precisely case 5.

Otherwise \(l<h\). For the lower boundary there are exactly
three possibilities. If \(l=a\), it is type 0; this point
cannot belong to \(R\), since the interval is open.
If \(a<l\) and \(l\in R\), it is a minimum because
it is a lower bound, and this is type 1. If \(a<l\) and
\(l\notin R\), it is type 2. The earlier inequality \(l<b\)
places the latter two boundaries strictly inside the outer interval.
These three cases are disjoint and exhaustive. The upper boundary
has the three disjoint exhaustive possibilities \(h=b\),
\(h<b\) and \(h\in R\), or \(h<b\) and \(h\notin R\).
They are respectively upper types 0, 1 and 2, with the internal
boundary strictly above \(a\).

There are nine ordered pairs of lower and upper types. In the source's
case numbering they are
\[
\begin{array}{c|ccccccccc}
 i&2&3&4&6&7&8&9&10&11\\ \hline
 (\lambda,\mu)&(0,0)&(0,1)&(0,2)&(1,0)&(1,1)&(1,2)&(2,0)&(2,1)&(2,2).
\end{array}
\]
Every pair retains \(l<h\). In particular the pair \((1,1)\)
here has two distinct attained extrema; it does not duplicate the
singleton case. Uniqueness of infimum and supremum gives unique
bounds, and their membership or nonmembership determines unique
types. Thus exactly one of the eleven cases holds.

\paragraph{The first-order boundary conditions are exact.}
We verify the formula descriptions so that the decomposition is
available as a first-order construction. All quantified names are
fresh; the fixed endpoint inside \(H\) is retained.
The lower type-0 formula is
\[
 \forall s\bigl(a<s\to
       \exists t\,(a<t<b\land t<s\land H(t))\bigr).
 \tag{1}
\]
Choosing \(s=b\) supplies a member of \(R\).
The point \(a\) is a lower bound, and (1) says every larger
point \(s\) has a member below it, so no such \(s\) is
a lower bound. Hence \(a=\inf R\). Conversely, if
\(a=\inf R\) and \(s>a\), then \(s\) is not a lower
bound, so a member of \(R\) lies below it, giving (1).

The lower type-1 formula is
\[
 a<l<b\land H(l)\land
                  \forall t\,(a<t<l\to\neg H(t)).
 \tag{2}
\]
It says \(l\in R\) and no member is below \(l\), precisely
the minimum condition. Conversely a minimum satisfies each conjunct.
The lower type-2 formula is the body's displayed formula. Its last
clause, instantiated with \(s=b\), supplies a member strictly
above \(l\). The other clauses exclude every member at or
below \(l\), within the outer interval, so \(l\) is a
lower bound absent from \(R\). Every \(s>l\) has a member
below it by the last clause, hence is not a lower bound. Thus
\(l\) is the infimum. Conversely an unattained interior infimum
has no member at or below it, and every larger point fails to be
a lower bound, giving every displayed conjunct.

The upper type-0 formula is
\[
 \forall s\bigl(s<b\to
       \exists t\,(a<t<b\land s<t\land H(t))\bigr).
 \tag{3}
\]
Its instance \(s=a\) supplies a member. The point \(b\)
is an upper bound; every smaller point fails to be one by (3),
so \(b=\sup R\). Conversely a supremum equal to \(b\)
gives such a member above every \(s<b\).
The upper type-1 formula is
\[
 a<h<b\land H(h)\land
                  \forall t\,(h<t<b\to\neg H(t)).
 \tag{4}
\]
It is equivalent to membership at \(h\) and absence of later
members, precisely the maximum condition. The upper type-2 formula
in the body excludes a member at or above \(h\), gives an
actual member through its instance \(s=a\), and makes every
smaller point fail to be an upper bound. It is therefore equivalent
to \(h\) being an unattained interior supremum. An unattained
supremum conversely supplies each of these properties, including a
member between any \(s<h\) and \(h\).

Emptiness is expressed by
\(\forall t\,(a<t<b\to\neg H(t))\).
The singleton descriptor at \(c\) is
\[
 a<c<b\land H(c)\land
 \forall t\bigl((a<t<c\lor c<t<b)\to\neg H(t)\bigr).
 \tag{5}
\]
Linearity partitions \((a,b)\) into its left subinterval,
the point \(c\), and its right subinterval, so (5) is
equivalent to \(R=\{c\}\). These formulas require no
object-language equality symbol.

The classification uses completeness solely to produce the bounded
set's infimum and supremum. The formula-to-bound equivalences themselves
hold in any linear order when the displayed conditions are satisfied.
For example, on \(\mathbb Q\) with \(a=0,b=2\), the
nonempty set \(\{q\in\mathbb Q:0<q<\sqrt2\}\) has no
supremum in the domain. It has no maximum and is not cofinal in
2, so none of the three upper types applies. This demonstrates
why the complete-order exhaustion is not asserted over all
\(\mathsf{Lin}\) \cite[ch.~III, p.~85]{Kamp1968}.
\end{proof}
\end{annexitem}

\begin{annexitem}{annex:p9.kamp.hull.interpolation}{thm:p09-kamp-hull-interpolation}{Both orientations of the required bracketing argument}
\begin{proof}
Let \(R=R_{a,b}\ne\varnothing\) have \(l=\inf R<h=\sup R\),
and take \(l<t<h\). If no member of \(R\) were strictly
below \(t\), then \(t\) would be a lower bound larger than
\(l\), contrary to the greatest-lower-bound property. Thus
there is \(r^-\in R\) with \(r^-<t\).
If no member were strictly above \(t\), then \(t\) would
be an upper bound smaller than \(h\), contrary to the
least-upper-bound property. Thus there is \(r^+\in R\)
with \(t<r^+\). Their membership in the original open interval gives
\[
                           a<r^-<t<r^+<b.
 \tag{1}
\]
This supplies actual witnesses on both sides, whether the bounds
are attained or unattained.

In the left orientation, \(H(s)=\phi(a,s)\). The two
members of \(R\) give \(\phi(a,r^-)\) and
\(\phi(a,r^+)\). Apply \cref{thm:p09-kamp-iii2-lemma6}
to the four strictly ordered points \(a,r^-,t,r^+\) and
this same fixed basic pattern. Its conclusion is precisely
\[
                   \phi(a,t)\longleftrightarrow
                               \exists r<t\;\phi(r,t).
\]
This is \(H(t)\longleftrightarrow P(t)\), with the
unrestricted strict-past witness required by the source.

For the right orientation we give the order reversal explicitly.
Write the fixed word as
\(W=(A_0,B_1,A_1,\ldots,B_n,A_n)\), and let
\[
                    W^{\mathrm{rev}}=(A_n,B_n,A_{n-1},\ldots,B_1,A_0).
\]
Keep the same domain and predicate interpretations, but use the
opposite strict order \(<^{\mathrm{op}}\), where
\(s<^{\mathrm{op}}t\) means \(t<s\) in the original order.
Construct a past basic pattern \(\phi^{\mathrm{rev}}\)
with word \(W^{\mathrm{rev}}\) in that opposite order.

For any original ordered pair \(s<q\), a realization
\(s=a_0<\cdots<a_n=q\) of \(W\) reverses to
\(q=a_n<^{\mathrm{op}}\cdots<^{\mathrm{op}}a_0=s\).
Each point retains its corresponding reversed point label, and
each opposite-order gap is the same set of points as the matching
original gap, with its corresponding reversed interval label.
This is a realization of \(W^{\mathrm{rev}}\). Reversing
again proves the converse. By the basic-word realization theorem,
\[
 \phi(s,q)\text{ in the original structure}
 \quad\Longleftrightarrow\quad
 \phi^{\mathrm{rev}}(q,s)\text{ in the opposite structure}.
 \tag{2}
\]
The argument is guarded by the ordered endpoints and preserves
the exact number of intervals.

Now \(H(r^-)\) and \(H(r^+)\) mean
\(\phi(r^-,b)\) and \(\phi(r^+,b)\). By (2) they
give the two fixed-start realizations
\(\phi^{\mathrm{rev}}(b,r^-)\) and
\(\phi^{\mathrm{rev}}(b,r^+)\). The four points
\(b,r^+,t,r^-\) are strictly increasing in the opposite
order, by (1). Lemma~6 in that linear order yields
\[
 \phi^{\mathrm{rev}}(b,t)
 \quad\Longleftrightarrow\quad
 \exists q<^{\mathrm{op}}t\;
                      \phi^{\mathrm{rev}}(q,t).
\]
On the left, (2) gives \(\phi(t,b)=H(t)\).
On the right, the guard means \(t<q\), and (2) gives
\(\phi(t,q)\). Thus the right side is precisely
\(\exists q>t\;\phi(t,q)=P(t)\). This proves the
right-orientation equivalence in full.

Only linearity and the existence of the stated bounds were used
to obtain (1). Lemma~6 applies on every linear order. Therefore
the interpolation has the conditional \(\mathsf{Lin}\)
scope asserted in the body; it does not assert that arbitrary
bounded truth sets possess extrema on incomplete orders.
\end{proof}
\end{annexitem}

\begin{annexitem}{annex:p9.kamp.corrected.boundary.decomposition}{thm:p09-kamp-corrected-boundary-decomposition}{Every guarded branch and the necessary fixed-endpoint conditions}
\begin{proof}
Fix the structure, its predicate interpretation, and its original
ordered endpoints \(a<b\). Every occurrence of these endpoints
inside the remaining disjunction \(D\) is held fixed.
Abbreviate the target by
\[
                    X=\forall t\,(a<t<b\to D(t)\lor H(t)).
\]
First work in any linear order under one of the stated descriptors.

Under \(\Gamma_1\), \(H(t)\) is false throughout
\((a,b)\), so \(X\) is equivalent to
\(\mathcal V_D(a,b)\). Consequently
\(\Gamma_1\land X\longleftrightarrow Z_1\).
Under the singleton descriptor \(\Gamma_5(c)\),
the point \(c\) satisfies \(H\) and every other point
of \((a,b)\) fails \(H\). The original universal condition
is therefore equivalent to \(D\) throughout \((a,c)\)
and throughout \((c,b)\), with no \(D\) requirement at
\(c\). Existentially choosing the singleton point gives
exactly \(Z_5\), in both directions.

For any two-boundary descriptor \(\Gamma_i(l,h)\), we have
the following complete partition of the tested interval:
\[
       (a,l),\quad \{l\}\text{ if }a<l,
       \quad(l,h),\quad \{h\}\text{ if }h<b,
       \quad(h,b).
 \tag{1}
\]
There is no overlap between its pieces, since \(l<h\),
and linearity places every \(a<t<b\) into one of them.
The first or last open interval can be empty, as can the middle
one on a discrete order.

Every point below \(l\) or above \(h\), within the original
interval, fails \(H\): the infimum is a lower bound and
the supremum an upper bound of the truth set. Thus \(D\)
is required throughout the two outer pieces of (1).
If the lower boundary is attained, it satisfies \(H\), so
no \(D(l)\) requirement remains. If it is an unattained
interior boundary, \(H(l)\) is false and \(D(l)\) is
required. If \(l=a\), the point is not in the original
open interval and no test is made there. The same three cases
hold at the upper boundary. These are precisely the tests
\(E_\lambda(l)\) and \(E_\mu(h)\) in the definition.

On the middle piece, \cref{thm:p09-kamp-hull-interpolation}
gives \(H(t)\longleftrightarrow P(t)\) at every point.
Therefore the universal condition there is exactly
\(\mathcal V_{D\lor P}(l,h)\). This proves both directions
of the branch equivalence under its descriptor: \(X\) supplies
every regional condition just listed; conversely those conditions
give \(D(t)\lor H(t)\) at each point in the exhaustive
partition (1).

For clarity, all nine regional matrices are the following. Each
matrix is conjoined with its own \(\Gamma_i\), and only the
interior bounds displayed in that row are existentially quantified:
\[
\begin{array}{c|l}
 2&\mathcal V_{D\lor P}(a,b)\\
 3&\mathcal V_{D\lor P}(a,h)\land\mathcal V_D(h,b)\\
 4&\mathcal V_{D\lor P}(a,h)\land D(h)\land\mathcal V_D(h,b)\\
 6&\mathcal V_D(a,l)\land\mathcal V_{D\lor P}(l,b)\\
 7&\mathcal V_D(a,l)\land\mathcal V_{D\lor P}(l,h)\land\mathcal V_D(h,b)\\
 8&\mathcal V_D(a,l)\land\mathcal V_{D\lor P}(l,h)\land D(h)
                                      \land\mathcal V_D(h,b)\\
 9&\mathcal V_D(a,l)\land D(l)\land\mathcal V_{D\lor P}(l,b)\\
 10&\mathcal V_D(a,l)\land D(l)\land\mathcal V_{D\lor P}(l,h)
                                      \land\mathcal V_D(h,b)\\
 11&\mathcal V_D(a,l)\land D(l)\land\mathcal V_{D\lor P}(l,h)
                                      \land D(h)\land\mathcal V_D(h,b).
\end{array}
\]
Rows 2--4 have the original initial endpoint as their lower bound;
rows 6--8 have an attained lower bound; rows 9--11 have an
unattained lower bound and therefore test \(D(l)\).
Rows 2, 6 and 9 reach the outer final endpoint. Rows 3, 7 and
10 have an attained upper bound, while rows 4, 8 and 11 have
an unattained upper bound and test \(D(h)\). This verifies
every combination, including both missing-boundary tests in row 11.

Now assume the order belongs to \(\mathsf{Com}\).
By \cref{thm:p09-kamp-eleven-boundaries}, the truth set has
exactly one of these descriptors, with its actual interior bounds
when present. If \(X\) holds, use those bounds and the
corresponding conditional branch equivalence to obtain \(Z_i\)
for that case. Hence the eleven-term disjunction holds.
Conversely, a true \(Z_i\) supplies its descriptor and any
required interior bounds; the conditional branch equivalence gives
\(X\). This direction does not require completeness.
The first direction uses it only for the existence of a descriptor.

All reasoning used the actual fixed endpoint in \(H\).
For a member \(\phi(x,z)\), introducing a lower bound \(l\)
does not replace the free \(x\) inside \(\phi(x,z)\) by
\(l\). It changes only the interval in which the original
truth set is examined. Likewise neither free endpoint inside
\(D\) is altered by dividing the interval. The opposite
orientation is covered by the fully proved right interpolation,
not by assuming that past and future formulas have identical syntax.

\paragraph{Verification of the finite bracketing counterexample.}
For the displayed one-interval pattern,
\[
 \phi(r,s)=A(r)\land\forall t\,(r<t<s\to B(t))\land C(s).
\]
At \((0,1)\), the initial and final predicates hold and the
gap is empty, so \(H(1)\) is true.
At \((0,4)\), both endpoint predicates hold but the gap
contains 2, which is not in \(B\); hence \(H(4)\) is false.
The other interior candidates 2 and 3 are outside \(C\),
so their \(H\) values are false as well. Thus
\(R_{0,5}=\{1\}\).
At 4 the remaining disjunction is \(Q(4)\), which is false,
so the original universal formula is false.

The unrestricted past closure at 4 is nevertheless true: choose
the initial point 3. We have \(A(3)\), \(C(4)\), and
the gap \((3,4)\) is empty, so \(\phi(3,4)\) holds.
Consider the weakened condition
\[
 \exists c\bigl(0<c<5\land H(c)\land
        \mathcal V_Q(0,c)\land\mathcal V_{Q\lor P}(c,5)\bigr).
 \tag{2}
\]
It holds with \(c=1\). The lower tested interval is empty.
In the upper interval the points 2 and 3 satisfy \(Q\),
and the point 4 satisfies \(P\), as just proved. This is
the formula shape printed as \(\chi_6\) on source p.~88.
The target is false, so its unconditional implication to the
target cannot be used.

Our descriptor for case 6 also requires the truth set to have
supremum 5, which it does not: its supremum is 1. Therefore
that branch is unavailable. The actual singleton branch requires
\(Q\) throughout \((1,5)\) and fails at 4, correctly
agreeing with the target. The finite order belongs to
\(\mathsf{Com}\), so assuming completeness does not repair
the weakened implication. The retained descriptors supply the
needed fixed-endpoint brackets rather than inferring them from
the existence of an unrelated past realization
\cite[ch.~III, pp.~85--89]{Kamp1968}.

This completes the universal boundary decomposition, including all
eleven cases. It does not establish elementary interval membership
for its descriptors or branches; that is the remaining substantive
step in the full Lemma~8 induction.
\end{proof}
\end{annexitem}
